\documentclass[aspectratio=169,11pt]{beamer}

\usetheme{Madrid}
\setbeamertemplate{navigation symbols}{}
\setbeamertemplate{footline}[frame number]

\usepackage[T1]{fontenc}
\usepackage[utf8]{inputenc}
\usepackage{lmodern}
\usepackage{microtype}
\usepackage{tikz}
\usetikzlibrary{positioning,calc}
\usepackage{booktabs}
\usepackage{tabularx}
\usepackage{array}
\usepackage{ragged2e}

\definecolor{Ink}{HTML}{1D2433}
\definecolor{Teal}{HTML}{1F6F78}
\definecolor{Gold}{HTML}{F2B134}
\definecolor{Sky}{HTML}{EAF6F6}
\definecolor{Mint}{HTML}{E7F3EC}
\definecolor{Rose}{HTML}{FCE9E9}
\definecolor{Sand}{HTML}{F8F4EC}

\setbeamercolor{palette primary}{bg=Ink,fg=white}
\setbeamercolor{palette secondary}{bg=Teal,fg=white}
\setbeamercolor{palette tertiary}{bg=Gold,fg=Ink}
\setbeamercolor{structure}{fg=Teal}
\setbeamercolor{frametitle}{bg=Sand,fg=Ink}
\setbeamercolor{block title}{bg=Teal,fg=white}
\setbeamercolor{block body}{bg=Sky,fg=Ink}
\setbeamercolor{alerted text}{fg=Gold}

\setbeamertemplate{itemize item}{\color{Gold}$\blacktriangleright$}
\setbeamertemplate{itemize subitem}{\color{Teal}$\circ$}

\newcommand{\claimbox}[1]{
  \begin{block}{Big idea}
  \Large #1
  \end{block}
}

\newcommand{\routinecard}[3]{
  \begin{tikzpicture}
    \node[
      draw=Teal,
      rounded corners=5pt,
      fill=#1,
      text width=0.28\linewidth,
      minimum height=2.6cm,
      inner sep=8pt,
      align=left
    ]{
      \textbf{\large #2}\par
      \vspace{0.15cm}
      #3
    };
  \end{tikzpicture}
}

\title{Impacting Practice}
\subtitle{From Access to Agency in K--5 Computer Science}
\author{Piter Garcia}
\institute{EDU 442 \\ Teaching Placement at Pine Brook IGNITE}
\date{Spring 2026}

\begin{document}

\begin{frame}
  \titlepage
  \vspace{0.15cm}
  {\small Built from lesson planning, placement reflection, and classroom patterns across kindergarten through grade 5.}
\end{frame}

\begin{frame}{Why this project?}
\claimbox{My placement kept showing me that engagement is not the same thing as access.}

\begin{columns}[T]
  \column{0.48\textwidth}
  \begin{block}{What I saw}
  \begin{itemize}
    \item Students were excited by robots, coding, Minecraft, LEGO, circuits, and design tools.
    \item Many lessons were hands-on, creative, and full of energy.
    \item Some students still got stuck before they could show what they knew.
  \end{itemize}
  \end{block}

  \column{0.48\textwidth}
  \begin{block}{What kept getting in the way}
  \begin{itemize}
    \item Hidden steps in directions
    \item Too much to hold in working memory
    \item Unclear help routines
    \item Adults becoming the only path forward
  \end{itemize}
  \end{block}
\end{columns}
\end{frame}

\begin{frame}{Placement context}
\begin{itemize}
  \item My teaching placement was at Pine Brook IGNITE across K--5 computer science and technology-rich lessons.
  \item The workspace documents show recurring lesson families: coding, robotics, digital design, Book Creator, Canva, Google Slides, Minecraft, and STEM challenges.
  \item The lesson packets repeatedly emphasize visual supports, chunked steps, flexible entry points, and clear proof of learning.
  \item Non-lesson reflection notes also highlight two important tensions:
  \begin{itemize}
    \item balancing support with independence
    \item treating care, regulation, and SEL as part of instruction
  \end{itemize}
\end{itemize}

\vspace{0.2cm}
\begin{alertblock}{Core takeaway}
Inclusive teaching is not just about adding support. It is about designing a path that students can actually enter, stay in, and own.
\end{alertblock}
\end{frame}

\begin{frame}{The artifact}
\begin{block}{My Impacting Practice artifact}
\Large The Access-to-Agency Routine
\end{block}

\vspace{0.2cm}
\begin{itemize}
  \item A short planning and facilitation routine for K--5 CS and digital learning lessons
  \item Designed to be ADHD-friendly, inclusion-minded, and realistic for busy classroom teaching
  \item Built from the repeated needs I saw in placement rather than from one single lesson
  \item Goal: reduce unnecessary dependence on adult rescue and increase student agency
\end{itemize}
\end{frame}

\begin{frame}{What the routine solves}
\begin{columns}[T]
  \column{0.49\textwidth}
  \begin{block}{Before}
  \begin{itemize}
    \item Students ask for help before trying
    \item ``I am done'' often means ``I am stuck''
    \item Strong ideas disappear under setup friction
    \item The teacher becomes the bottleneck
  \end{itemize}
  \end{block}

  \column{0.49\textwidth}
  \begin{block}{After}
  \begin{itemize}
    \item Students know the first move
    \item Students can re-enter the task after confusion
    \item Help becomes structured instead of chaotic
    \item Success is visible and explainable
  \end{itemize}
  \end{block}
\end{columns}
\end{frame}

\begin{frame}{The routine in five moves}
\centering
\routinecard{Mint}{1. See}{Model the task and name one visible goal.}
\hspace{0.02\linewidth}
\routinecard{Sky}{2. Start}{Give only the first one or two moves.}
\hspace{0.02\linewidth}
\routinecard{Mint}{3. Try}{Students work with a visible help path.}

\vspace{0.25cm}
\routinecard{Sky}{4. Check}{Pause once to reset before confusion spreads.}
\hspace{0.02\linewidth}
\routinecard{Rose}{5. Show}{Students explain, save, test, or share proof.}
\end{frame}

\begin{frame}{Why this is ADHD-friendly}
\begin{columns}[T]
  \column{0.48\textwidth}
  \begin{block}{It lowers cognitive load}
  \begin{itemize}
    \item One goal at a time
    \item Fewer hidden steps
    \item Repeated lesson structure
    \item Visual anchors students can re-find
  \end{itemize}
  \end{block}

  \column{0.48\textwidth}
  \begin{block}{It protects momentum}
  \begin{itemize}
    \item Short chunks of attention
    \item Quick resets instead of long derailments
    \item Movement and hands-on learning with purpose
    \item Structured choice without overload
  \end{itemize}
  \end{block}
\end{columns}

\vspace{0.2cm}
\claimbox{The goal is not perfect focus. The goal is a lesson structure students can come back to.}
\end{frame}

\begin{frame}{What I learned from placement materials}
\small
\begin{tabularx}{\textwidth}{>{\RaggedRight\arraybackslash}p{0.31\textwidth} >{\RaggedRight\arraybackslash}X}
\toprule
\textbf{Repeated pattern} & \textbf{What it taught me} \\
\midrule
Visual supports and pictures & Access improves when directions are concrete, not only verbal. \\
Chunked tasks and short signals & Students need structure that matches real attention patterns. \\
Choice boards and flexible responses & Entry points matter; students do not all need the same doorway into learning. \\
Debug and reset moments & Productive struggle needs a map, not just encouragement. \\
Proof of learning prompts & Students need to know what counts as success before frustration takes over. \\
\bottomrule
\end{tabularx}
\end{frame}

\begin{frame}{Examples from my placement}
\begin{itemize}
  \item In Grade 1 lesson packets, the supports repeatedly stress simple language, lots of pictures, short chunks, quiet breaks, and visible checklists.
  \item In Grade 4 audience notes, the recurring need is clear setup routines, visible success criteria, and prompts that support debugging without rescue.
  \item Across multiple lesson packets, one revision move appears again and again: make the next step, help routine, and success condition more visible.
  \item Across the non-lesson context notes, care and regulation are treated as real parts of readiness, not side issues.
\end{itemize}

\vspace{0.15cm}
\begin{alertblock}{My synthesis}
The access problem was not a lack of student ability. It was often a mismatch between task demands and the supports built into the lesson structure.
\end{alertblock}
\end{frame}

\begin{frame}{How EDU 442 shaped this project}
\begin{itemize}
  \item \textbf{Identity-conscious teaching:} students need pathways into learning that do not require them to mask who they are.
  \item \textbf{Culturally sustaining pedagogy:} language, collaboration, creativity, and student ways of making meaning are resources, not distractions.
  \item \textbf{Disability and neurodivergence lenses:} access needs should not be misread as laziness, defiance, or inability.
  \item \textbf{Joy and refusal:} engagement matters, but joy is strongest when students have some power, not just compliance.
  \item \textbf{Structural thinking:} when many students get stuck, the first question should be about lesson design, not about student deficits.
\end{itemize}
\end{frame}

\begin{frame}{What this looks like in practice}
\begin{columns}[T]
  \column{0.5\textwidth}
  \begin{block}{Teacher planning moves}
  \begin{itemize}
    \item Write one student-friendly goal
    \item Identify the first two steps
    \item Build one visible help routine
    \item Plan one midpoint reset
    \item Decide what proof students will show
  \end{itemize}
  \end{block}

  \column{0.46\textwidth}
  \begin{block}{Student-facing experience}
  \begin{itemize}
    \item I know what I am doing
    \item I know how to begin
    \item I know what to try if I get stuck
    \item I know what success looks like
    \item I can show what I learned
  \end{itemize}
  \end{block}
\end{columns}
\end{frame}

\begin{frame}{A sample classroom translation}
\begin{block}{Example: robotics, coding, or digital creation lesson}
\begin{enumerate}
  \item Show a model and name the goal in plain language.
  \item Put the first two moves on the screen.
  \item Give students a visible help path: partner, checklist, anchor image, then teacher.
  \item Pause once to debug as a whole group before frustration becomes shutdown.
  \item End with one proof move: save, test, explain, screenshot, or demonstrate.
  \end{enumerate}
\end{block}

\vspace{0.15cm}
This routine works across lesson types because it is not tied to one app or one grade. It is a structure for access.
\end{frame}

\begin{frame}{What changed in my thinking}
\begin{columns}[T]
  \column{0.48\textwidth}
  \begin{block}{At first}
  \begin{itemize}
    \item I often thought more support meant better support.
    \item I focused on making lessons engaging and creative.
    \item I noticed breakdowns but did not always have a framework for them.
  \end{itemize}
  \end{block}

  \column{0.48\textwidth}
  \begin{block}{Now}
  \begin{itemize}
    \item I think more about access, entry points, and independence.
    \item I design for re-entry, not just for attention.
    \item I see lesson structure itself as an equity practice.
  \end{itemize}
  \end{block}
\end{columns}
\end{frame}

\begin{frame}{Limits and next steps}
\begin{itemize}
  \item This project is a practical routine, not a full solution to institutional inequity.
  \item It does not remove the need for staffing, planning time, family partnership, or teacher voice.
  \item It works best when paired with classroom knowledge, flexibility, and relational care.
  \item My next step would be to pilot the routine across multiple lessons and collect student work, teacher notes, and reflection data on independence and participation.
\end{itemize}
\end{frame}

\begin{frame}{Closing}
\claimbox{My impacting practice claim is simple: when access is built into the structure of a lesson, students are more likely to stay in the work, make choices, and show agency.}

\vspace{0.2cm}
\begin{itemize}
  \item This project grows directly out of my teaching placement.
  \item It translates EDU 442 ideas into a usable classroom tool.
  \item It helps me teach in a way that is more inclusive, more humane, and more realistic.
\end{itemize}
\end{frame}

\end{document}
